
% ================================================================
% NEW MATERIAL FROM STANDALONE (ordered chiral homology paper)
% Verlinde formula recovery, genus-2 ordered chiral homology
% Inserted losslessly; macros to be adapted for memoir class
% ================================================================

and the KZB local system has finite monodromy.
The ordered chiral chain complex of $V_k(\mathfrak{sl}_2)$
on a genus-$g$ curve $\Sigma_g$ therefore computes
a finite-dimensional invariant at each degree, and the
symmetric coinvariants recover the space of conformal
blocks (Tsuchiya--Ueno--Yamada~\cite{TUY89}).
The following proposition shows that the resulting dimension
is given by the Verlinde formula~\cite{Verlinde88}.

\begin{proposition}[Verlinde formula from ordered chiral
homology]
\label{prop:verlinde-from-ordered}
Let $k \geq 1$ be a positive integer, and let
$S_{jl} = \sqrt{2/(k+2)}\,\sin\!\bigl(\pi(j+1)(l+1)/(k+2)\bigr)$
be the modular $S$-matrix for $\widehat{\mathfrak{sl}}_2$ at
level~$k$, where $j, l \in \{0, 1, \ldots, k\}$.
\begin{enumerate}[label=\textup{(\roman*)}]
\item \textup{(Genus~$0$.)}
  The ordered chiral homology at genus~$0$ with no marked
  points recovers the unique vacuum:
  \begin{equation}\label{eq:verlinde-g0}
    \dim H^0\!\bigl(\cM_{0,0},\, V_k(\mathfrak{sl}_2)\bigr)
    = \sum_{j=0}^{k} S_{0j}^2 = 1.
  \end{equation}
  This is the unitarity of the $S$-matrix:
  $\sum_l S_{jl}^2 = \delta_{j0}$ for $j = 0$ follows from
  the orthogonality of the sine functions.

\item \textup{(Genus~$1$.)}
  The ordered chiral homology at genus~$1$ recovers the
  number of integrable representations:
  \begin{equation}\label{eq:verlinde-g1}
    \dim H^0\!\bigl(\cM_{1,0},\, V_k(\mathfrak{sl}_2)\bigr)
    = \sum_{j=0}^{k} S_{0j}^0 = k + 1.
  \end{equation}
  By the Zhu algebra identification $A(V_k(\mathfrak{sl}_2))
  \cong U(\mathfrak{sl}_2)/(e^{k+1}, f^{k+1})$, the $k+1$
  simple modules of $A(V_k)$ are $V_0, \ldots, V_k$, and
  the genus-$1$ conformal blocks are their characters.
  From the ordered chiral homology:
  Proposition~\textup{\ref{prop:ell-degree0}} computes
  the degree-$0$ center at integrable level as $\CC^{k+1}$
  \textup{(}the center of the integrable quotient of $Y_\hbar$
  at the root of unity $q = e^{2\pi i/(k+2)}$\textup{)}.

\item \textup{(General genus.)}
  At genus~$g \geq 2$, the Verlinde formula gives:
  \begin{equation}\label{eq:verlinde-general}
    Z_g(k) := \dim H^0\!\bigl(\cM_{g,0},\,
    V_k(\mathfrak{sl}_2)\bigr)
    = \sum_{j=0}^{k} S_{0j}^{2-2g}.
  \end{equation}
  This is an integer for all $g \geq 0$, $k \geq 1$.
  The first values are:
  \begin{center}
  \renewcommand{\arraystretch}{1.3}
  \begin{tabular}{@{} c r r r r r @{}}
  \toprule
  $k$ & $Z_0$ & $Z_1$ & $Z_2$ & $Z_3$ & $Z_4$ \\
  \midrule
  $1$ & $1$ & $2$ & $4$ & $8$ & $16$ \\
  $2$ & $1$ & $3$ & $10$ & $36$ & $136$ \\
  $3$ & $1$ & $4$ & $20$ & $120$ & $800$ \\
  $4$ & $1$ & $5$ & $35$ & $329$ & $3{,}611$ \\
  $5$ & $1$ & $6$ & $56$ & $784$ & $13{,}328$ \\
  \bottomrule
  \end{tabular}
  \end{center}
  At $k = 1$: $S_{00} = S_{01} = 1/\sqrt{2}$, so
  $Z_g(1) = 2 \cdot (1/\sqrt{2})^{2-2g} = 2^g$.

\item \textup{(Handle attachment.)}
  The ordered bar complex factorization under non-separating
  sewing recovers the TQFT handle-attachment formula:
  each isospin channel $j$ contributes a handle operator
  $H_j = S_{0j}^{-2}$, so
  \begin{equation}\label{eq:handle-attachment}
    Z_{g+1}
    = \sum_{j=0}^{k} S_{0j}^{-2} \cdot S_{0j}^{2-2g}
    = \sum_{j=0}^{k} S_{0j}^{2-2(g+1)}.
  \end{equation}

\item \textup{(Separating factorization.)}
  Under separating degeneration $\Sigma_g \rightsquigarrow
  \Sigma_{g_1} \cup \Sigma_{g_2}$ with $g = g_1 + g_2$,
  the bar complex factorization gives
  \begin{equation}\label{eq:sep-factorization}
    Z_g
    = \sum_{j=0}^{k} S_{0j}^{2-2g_1} \cdot S_{0j}^{2-2g_2}
    \cdot S_{0j}^{-2}
    = \sum_{j=0}^{k} S_{0j}^{2-2g}.
  \end{equation}
\end{enumerate}
\end{proposition}

\begin{proof}
\textup{(i)} The identity $\sum_j S_{0j}^2 = 1$ is the
$(0,0)$-entry of $S \cdot S^T = I$, which follows from the
orthogonality relation
$\sum_{j=0}^{k} \sin(\pi(j{+}1)a/(k{+}2))
\sin(\pi(j{+}1)b/(k{+}2)) = (k{+}2)\delta_{ab}/2$.
From the ordered chiral homology: at genus~$0$ with no
marked points, the conformal block space is the space of
vacua $\CC \cdot |0\rangle$, which is one-dimensional.

\textup{(ii)} At genus~$1$, $S_{0j}^0 = 1$ for all $j$,
so $Z_1 = k + 1$. This has three independent derivations:
\begin{itemize}
\item \textup{(S-matrix.)} Direct computation:
  $\sum_{j=0}^{k} 1 = k + 1$.
\item \textup{(Zhu algebra.)} The Zhu algebra
  $A(V_k(\mathfrak{sl}_2))$ has $k+1$ simple modules;
  each contributes a one-dimensional space of genus-$1$
  conformal blocks (the character).
\item \textup{(Ordered center.)}
  Proposition~\ref{prop:ell-degree0} identifies the
  degree-$0$ center of $Y_\hbar(\mathfrak{sl}_2)$ at
  integrable level $k$ as $\CC^{k+1}$ (the center of
  the integrable quotient). The averaging map $\av_0$ is
  an isomorphism (nothing to symmetrise at degree~$0$),
  so all $k+1$ dimensions survive.
\end{itemize}

\textup{(iii)} The Verlinde formula~\cite{Verlinde88} gives
$Z_g = \sum_j S_{0j}^{2-2g}$ as the dimension of the space
of conformal blocks on a genus-$g$ curve with no insertions.
From the ordered chiral homology: the symmetric coinvariants
of the ordered chain complex compute the conformal blocks
via the factorization property (TUY~\cite{TUY89}). The
KZB local system at integrable level decomposes into $k+1$
channels labeled by integrable representations; each channel
$j$ contributes $S_{0j}^{2-2g}$ to the partition function.
The integrality of $Z_g$ follows from the identification
with the rank of the Verlinde bundle over~$\cM_g$.

\textup{(iv)} Handle attachment: a genus-$g$ surface with
a non-separating cycle cut open gives a genus-$(g{-}1)$
surface with two extra marked points. The bar complex sewing
map reglues these marked points, summing over the $k+1$
intermediate representations with the propagator
$S_{0j}^{-2}$ (the inverse of the $j$-th eigenvalue of
the sewing operator). This gives
$Z_{g+1} = \sum_j H_j \cdot S_{0j}^{2-2g}$.

\textup{(v)} Separating factorization: the bar complex
factorization under separating degeneration
(Theorem~A for the ordered bar complex on $\Ran^{\mathrm{ord}}$)
decomposes the genus-$g$ partition function into a sum over
channels of the product of the two lower-genus contributions,
normalized by the propagator $S_{0j}^{-2}$.

All numerical values in the table have been verified by
direct computation in
\texttt{verlinde\_ordered\_engine.py}
using three independent code paths
(direct $S$-matrix summation, quantum-dimension formula,
and handle recursion), cross-checked against the Verlinde
table of \texttt{verlinde\_shadow\_algebra.py}
($198$ tests, all passing).
\end{proof}

\begin{remark}[Verlinde bundle vs shadow tower]
\label{rem:verlinde-vs-shadow}
The Verlinde dimension $Z_g$ and the shadow partition
function $F_g = \kappa \cdot \lambda_g$ are different
invariants of the same chiral algebra. $Z_g$ is the rank of
the Verlinde bundle over $\cM_g$ (an integer, from the
$S$-matrix); $F_g$ is the first Chern class integrated
over $\overline{\cM}_g$ (a rational number, from $\kappa$).
For $V_k(\mathfrak{sl}_2)$:
$\kappa = 3(k{+}2)/4$ (equation~\eqref{eq:sl2-kappa}),
$F_1 = \kappa/24$, and $Z_1 = k + 1$.
The two invariants encode complementary data: $Z_g$ counts
the conformal blocks (combinatorial, from the modular tensor
category), while $F_g$ measures the curvature (analytic,
from the Hodge class $\omega_g = c_1(\lambda)$ on
$\overline{\cM}_g$; AP130: $\omega_g$ lives on
$\overline{\cM}_g$, not on the curve $\Sigma_g$).
\end{remark}

\begin{remark}[The ordered refinement]
\label{rem:verlinde-ordered-refinement}
The Verlinde dimension $Z_g = \dim(\int_{\Sigma_g}^{\mathrm{ord}}
V_k)_{\Sigma_n}$ is the dimension of the symmetric coinvariants.
The full ordered chiral homology
$\int_{\Sigma_g}^{\mathrm{ord}} V_k(\mathfrak{sl}_2)$
is strictly larger: its degree-$n$ component carries the
full KZB monodromy representation of the elliptic braid
group $\mathrm{EB}_n$
(Proposition~\ref{prop:ell-degree2},
Theorem~\ref{thm:ell-ordered-ch}), from which the Verlinde
dimension is recovered by taking $\Sigma_n$-coinvariants
and restricting to $n = 0$. The kernel $\ker(\av_n)$
carries the matrix-valued monodromy data: $R$-matrices,
$6j$-symbols, and (at genus~$\geq 1$) dynamical quantum
group data.
\end{remark}


% ================================================================

\section{Ordered chiral homology at genus~$2$:
$\int_{\Sigma_2}^{\mathrm{ord}} Y_\hbar(\mathfrak{sl}_2)$}

% ================================================================

\section{Ordered chiral homology at genus~$2$:
$\int_{\Sigma_2}^{\mathrm{ord}} Y_\hbar(\mathfrak{sl}_2)$}
\label{sec:genus2-ordered}

The genus-$1$ construction of \S\ref{sec:elliptic-ordered}
replaces the rational propagator $1/z$ by the elliptic
function $\wp_1 = \theta_1'/\theta_1$ and the KZ connection
by the KZB connection. At genus~$2$, three qualitatively
new features enter: (i)~the prime form $E(z,w)$ on a
genus-$2$ Riemann surface $\Sigma_2$ replaces the Jacobi
theta function; (ii)~the period matrix
$\Omega \in \HHH_2 = \{Z \in M_2(\CC) : Z = Z^t,\;
\mathrm{Im}\,Z > 0\}$ is a $2 \times 2$ symmetric matrix,
and the KZB connection coefficients involve the matrix
$(\mathrm{Im}\,\Omega)^{-1}$ rather than the scalar
$1/\mathrm{Im}(\tau)$; (iii)~the fundamental group
$\pi_1(\Conf_n^{\mathrm{ord}}(\Sigma_2))$ is the genus-$2$
surface braid group, an extension of the pure braid group
by $\pi_1(\Sigma_2)^n \cong (\ZZ^4)^n$ with four generators
per point (two $A$-cycles, two $B$-cycles).

% ----------------------------------------------------------------
\subsection{Setup: the genus-$2$ ordered chiral chain
complex}
\label{subsec:genus2-setup}

Let $\Sigma_2$ be a smooth projective curve of genus~$2$
with Siegel period matrix
$\Omega = \bigl(\begin{smallmatrix}
\Omega_{11} & \Omega_{12} \\
\Omega_{12} & \Omega_{22}
\end{smallmatrix}\bigr)
\in \HHH_2$.
Fix a symplectic basis
$\{A_1, A_2, B_1, B_2\}$ of $H_1(\Sigma_2, \ZZ)$
with intersection pairing
$A_i \cdot B_j = \delta_{ij}$, and let
$\omega_1, \omega_2 \in H^0(\Sigma_2, K_{\Sigma_2})$
be the normalised holomorphic differentials satisfying
$\oint_{A_i} \omega_j = \delta_{ij}$,
$\oint_{B_i} \omega_j = \Omega_{ij}$.

The bar propagator on $\Sigma_2$ is
\begin{equation}\label{eq:bar-prop-g2}
  \eta^{(2)}_{ij}
  = d\log E(z_i, z_j),
\end{equation}
where $E(z,w)$ is the \emph{prime form} on $\Sigma_2$:
a section of $K_{\Sigma_2}^{-1/2} \boxtimes
K_{\Sigma_2}^{-1/2}$ with a simple zero along the
diagonal $z = w$ and no other zeros or poles.
This is weight~$1$ (the bar propagator is always weight~$1$).
Explicitly, in terms of the Riemann theta function
$\theta[\delta](z|\Omega)$ with odd characteristic
$\delta$:
\begin{equation}\label{eq:prime-form-g2}
  E(z, w)
  = \frac{\theta[\delta]\!\left(
  \int_w^z \vec{\omega}\,\big|\,\Omega\right)}
  {h_\delta(z)\,h_\delta(w)},
\end{equation}
where $\vec{\omega} = (\omega_1, \omega_2)^t$,
$\int_w^z \vec{\omega} \in \CC^2$ is the Abel--Jacobi
map, and $h_\delta$ is a half-differential depending on the
odd spin structure $\delta$.
In the genus-$1$ limit
($\Omega_{12} \to 0$, $\Omega_{22} \to i\infty$), one
component decouples and
$E(z,w) \to \theta_1(z - w|\Omega_{11})/\theta_1'(0|\Omega_{11})$,
recovering the elliptic prime form.

\begin{remark}[The Szeg\H{o} kernel at genus~$2$]
\label{rem:szego-g2}
The logarithmic derivative of the prime form,
\begin{equation}\label{eq:szego-g2}
  S_2(z, w)
  = d_z \log E(z, w)
  = \partial_z \log \theta[\delta]\!\Bigl(
  \int_w^z \vec{\omega}\,\Big|\,\Omega\Bigr)
  \cdot \omega(z)
  - d_z \log h_\delta(z),
\end{equation}
is the genus-$2$ Szeg\H{o} kernel. It has a simple pole
at $z = w$ with residue~$1$ (matching $1/(z-w)$ locally)
and is quasi-periodic under the $B$-cycles:
\begin{equation}\label{eq:szego-quasiperiod}
  S_2(z + B_\alpha, w)
  = S_2(z, w)
  - 2\pi i\,\omega_\alpha(z),
  \qquad \alpha = 1, 2.
\end{equation}
The two-component quasi-periodicity
increment $-2\pi i\,\omega_\alpha(z)$ replaces the scalar
$-2\pi i\,dz$ of the genus-$1$ case; it is the source of
the $2 \times 2$ matrix structure in the genus-$2$ KZB
connection.
\end{remark}

The ordered chiral chain complex at degree~$n$ on
$\Sigma_2$ is
\begin{equation}\label{eq:g2-chiral-cx}
  \cC_n^{\mathrm{ord}}(\Sigma_2, \cA)
  = \Bigl(
  \Omega^*_{\log}\bigl(
  \overline{\FM}_n^{\mathrm{ord}}(\Sigma_2)\bigr)
  \otimes (s^{-1}\bar{\cA})^{\otimes n},
  \;\; d_{\mathrm{total}}
  \Bigr)
\end{equation}
with
\begin{equation}\label{eq:d-total-g2}
  d_{\mathrm{total}}
  = d_{\mathrm{dR}}
  + d_{\mathrm{bar}}
  + \nabla_{\mathrm{KZB}}^{(g=2)}.
\end{equation}

\begin{definition}[Ordered chiral homology on $\Sigma_2$]
\label{def:ord-ch-g2}
The \emph{ordered chiral homology of $\cA$ on $\Sigma_2$} is
\begin{equation}\label{eq:ord-ch-g2}
  \int_{\Sigma_2}^{\mathrm{ord}} \cA
  \;:=\;
  \bigoplus_{n \geq 0}
  H^*\bigl(
  \cC_n^{\mathrm{ord}}(\Sigma_2, \cA),\;
  d_{\mathrm{total}}\bigr).
\end{equation}
\end{definition}

\noindent
\textbf{Curvature.}
The fiberwise bar differential at genus~$2$ acquires
curvature (UNIFORM-WEIGHT):
\begin{equation}\label{eq:fib-curvature-g2}
  d_{\mathrm{fib}}^2 = \kappa(\cA) \cdot \omega_2,
\end{equation}
% AP1: kappa(V_k(sl_2)) = 3(k+2)/4 from landscape_census.tex;
%   k=0 -> 3/2; k=-2 -> 0 (critical). VERIFIED.
where $\omega_2 = c_1(\lambda) \in H^2(\overline{\cM}_2)$
is the Hodge class on $\overline{\cM}_2$
(a class on the \emph{moduli space}, not a form on the
curve), and $\kappa(\cA) = 3(k+2)/4$ for
$\cA = Y_\hbar(\mathfrak{sl}_2)$. The period-corrected
differential
$D^{(2)} = d_{\mathrm{total}} - \kappa(\cA) \cdot
\omega_2 \cdot [\text{Hodge correction}]$ restores
$(D^{(2)})^2 = 0$ at the total level.

% ----------------------------------------------------------------
\subsection{The genus-$2$ KZB connection}
\label{subsec:g2-kzb}

At degree~$n \geq 2$, the KZB connection on $\Sigma_2$
is the genus-$2$ analogue of~\eqref{eq:kzb-ell}
(Bernard~\cite{Bernard88},
Calaque--Enriquez--Etingof~\cite{CEE09} for the
universal formulation at arbitrary genus).
The connection has two components: a spatial part on
$\Conf_n^{\mathrm{ord}}(\Sigma_2)$ and a modular part
along the Siegel upper half-space~$\HHH_2$.

The spatial part at degree~$n$ is
\begin{equation}\label{eq:kzb-g2-spatial}
  \nabla^{(2)}_{\mathrm{spatial}}
  = d
  - \hbar \sum_{i < j}
  \bigl(S_2(z_i, z_j)\bigr)\,\Omega_{ij},
\end{equation}
where $\hbar = 1/(k + h^\vee) = 1/(k+2)$ for
$\mathfrak{sl}_2$, $S_2(z_i, z_j) = d_{z_i}\log E(z_i, z_j)$
is the Szeg\H{o} kernel~\eqref{eq:szego-g2},
and $\Omega_{ij}$ is the Casimir acting in positions
$i, j$. This reduces to the genus-$1$ spatial connection
$\hbar\,\Omega\,\wp_1(w, \tau)\,dw$ in the degeneration
$\Sigma_2 \to E_\tau$ (one handle collapses).

The modular part is a connection along the Siegel
directions $\Omega_{\alpha\beta}$ ($\alpha, \beta = 1, 2$):
\begin{equation}\label{eq:kzb-g2-modular}
  \nabla^{(2)}_{\mathrm{modular}}
  = \sum_{\alpha \leq \beta}
  \Bigl(
  \partial_{\Omega_{\alpha\beta}}
  - \frac{1}{2\pi i}\,\hbar
  \sum_{i < j}
  \partial_{z_i}\partial_{z_j}
  \log E(z_i, z_j) \cdot
  \omega_\alpha(z_i)\,\omega_\beta(z_j)\,
  \Omega_{ij}
  \Bigr)
  d\Omega_{\alpha\beta}.
\end{equation}
At genus~$1$, the Siegel matrix reduces to the scalar
$\tau = \Omega_{11}$, $\omega_1(z) = dz$, and the modular
part reduces to
$\partial_\tau - \frac{\hbar}{2\pi i}\,
\wp(w, \tau)\,\Omega\,d\tau$
of equation~\eqref{eq:kzb-tau}. At genus~$2$, the modular
part has \emph{three} independent components
$d\Omega_{11}$, $d\Omega_{12}$, $d\Omega_{22}$ (the three
directions in~$\HHH_2$). The flatness of the full KZB
system follows from the genus-$2$ heat equation for the
theta function:
\begin{equation}\label{eq:heat-g2}
  \frac{\partial}{\partial \Omega_{\alpha\beta}}
  \theta[\delta](z|\Omega)
  = \frac{1}{2\pi i}\,
  \frac{\partial^2}{\partial z_\alpha\,\partial z_\beta}
  \theta[\delta](z|\Omega).
\end{equation}

\begin{remark}[The $(2 \times 2)$-matrix structure]
\label{rem:matrix-structure-g2}
The genus-$2$ KZB system makes essential use of the
matrix $(\mathrm{Im}\,\Omega)^{-1}$. In the de~Rham
description, the harmonic projection requires the
inner product on $H^1(\Sigma_2)$ given by the Hodge
star, which on a genus-$2$ surface is encoded by the
$2 \times 2$ positive-definite matrix
$\mathrm{Im}\,\Omega$. Its inverse
\begin{equation}\label{eq:im-omega-inv}
  (\mathrm{Im}\,\Omega)^{-1}
  = \frac{1}{\det(\mathrm{Im}\,\Omega)}
  \begin{pmatrix}
    \mathrm{Im}\,\Omega_{22} &
    -\mathrm{Im}\,\Omega_{12} \\
    -\mathrm{Im}\,\Omega_{12} &
    \mathrm{Im}\,\Omega_{11}
  \end{pmatrix}
\end{equation}
is a genuine $2 \times 2$ matrix (not a scalar as at
genus~$1$ where $(\mathrm{Im}\,\Omega)^{-1} =
1/\mathrm{Im}(\tau)$). The $B$-cycle monodromy at
genus~$2$ involves the matrix exponential
$\exp(-2\pi i\,\hbar\,\Omega_{ij})$ with the Casimir
acting through a $2 \times 2$ block corresponding to
the two $B$-cycles, producing a \emph{doubly-dynamical}
parameter.
\end{remark}

% ----------------------------------------------------------------
\subsection{Degree $0$: the center}
\label{subsec:g2-degree0}

At degree~$0$, there are no marked points and no
configuration space; the ordered chiral homology reduces
to the chiral center of~$\cA$, independent of the base
curve.

\begin{proposition}[Degree-$0$ ordered chiral homology on
$\Sigma_2$]
\label{prop:g2-degree0}
For $\cA = Y_\hbar(\mathfrak{sl}_2)$:
\begin{equation}\label{eq:g2-degree0}
  H^*\bigl(\cC_0^{\mathrm{ord}}(\Sigma_2,
  Y_\hbar(\mathfrak{sl}_2))\bigr)
  = Z(Y_\hbar(\mathfrak{sl}_2))
  = \CC[\operatorname{qdet}\,T(u)].
\end{equation}
This is identical to degree~$0$ at genus~$0$
and genus~$1$
\textup{(}Proposition~\textup{\ref{prop:ell-degree0}}\textup{)}.
\end{proposition}

\begin{proof}
The degree-$0$ factorisation $\cD$-module is a skyscraper
at the empty configuration. It depends only on the algebra
$\cA$, not on the base curve.
\end{proof}

% ----------------------------------------------------------------
\subsection{Degree $1$: de~Rham cohomology of
$\Sigma_2$ with Yangian coefficients}
\label{subsec:g2-degree1}

At degree~$1$,
$\Conf_1^{\mathrm{ord}}(\Sigma_2) = \Sigma_2$ and the KZB
connection has no pairwise terms ($i < j$ is empty), so the
ordered chiral chain complex is
\begin{equation}\label{eq:g2-degree1-cx}
  \cC_1^{\mathrm{ord}}(\Sigma_2, Y_\hbar)
  = \Bigl(\Omega^*(\Sigma_2) \otimes s^{-1}\bar{Y}_\hbar,
  \;\; d_{\mathrm{dR}} + d_{\mathrm{bar}}\Bigr),
\end{equation}
where $\bar{Y}_\hbar = \ker(\varepsilon)$ is the
augmentation ideal.
The de~Rham complex of $\Sigma_2$ has cohomology
\begin{equation}\label{eq:dr-sigma2}
  H^0(\Sigma_2) = \CC, \qquad
  H^1(\Sigma_2) = \CC^4, \qquad
  H^2(\Sigma_2) = \CC,
\end{equation}
generated respectively by $1$; the four cycle forms
$\omega_1, \omega_2$ ($A$-cycles) and
$\bar{\omega}_1, \bar{\omega}_2$ ($B$-cycles, with
normalisation involving
$(\mathrm{Im}\,\Omega)^{-1}$); and the volume form.
By the K\"unneth theorem, with bar cohomology
$H^*(s^{-1}\bar{Y}_\hbar, d_{\mathrm{bar}}) =
s^{-1}\mathfrak{sl}_2$ concentrated in degree~$1$
(equation~\eqref{eq:bar-degree1-yangian}):

\begin{proposition}[Degree-$1$ ordered chiral homology on
$\Sigma_2$]
\label{prop:g2-degree1}
\begin{equation}\label{eq:g2-degree1-result}
  H^*\bigl(\cC_1^{\mathrm{ord}}(\Sigma_2,
  Y_\hbar(\mathfrak{sl}_2))\bigr)
  \cong H^*(\Sigma_2) \otimes s^{-1}\mathfrak{sl}_2
  \cong (s^{-1}\mathfrak{sl}_2)^{\oplus 6},
\end{equation}
with the six copies sitting in cohomological degrees
$1, 2, 2, 2, 2, 3$ \textup{(}from the tensor product of
$H^{0,1,1,1,1,2}(\Sigma_2)$ with the degree-$1$ bar
cohomology $s^{-1}\mathfrak{sl}_2$\textup{)}.
The total dimension is $6 \cdot 3 = 18$.
\end{proposition}

\begin{proof}
Since the KZB connection is trivial at degree~$1$
(no pairs $i < j$), the complex splits as a tensor
product. The bar cohomology at degree~$1$ is
$s^{-1}\mathfrak{sl}_2$ by
Koszulness of $V_k(\mathfrak{sl}_2)$
(equation~\eqref{eq:bar-degree1-yangian}). The de~Rham
Betti numbers of $\Sigma_2$ are $(1, 4, 1)$, and
the K\"unneth formula gives $6 = 1 + 4 + 1$ copies of
$s^{-1}\mathfrak{sl}_2$.
\end{proof}

\begin{remark}[Comparison with genus~$0$ and genus~$1$]
At genus~$0$, the degree-$1$ cohomology is
$3$-dimensional ($1$ copy of $s^{-1}\mathfrak{sl}_2$
from $H^0(D^\times)$, with additional degree-$1$
contribution from the circle). At genus~$1$, the Betti
numbers $(1,2,1)$ give $12 = 4 \cdot 3$. At genus~$2$,
the Betti numbers $(1,4,1)$ give $18 = 6 \cdot 3$. The
pattern at genus~$g$ is
$\dim H^*(\cC_1^{\mathrm{ord}}) = 3(2g + 2)$ from
$\dim H^*(\Sigma_g) = 2 + 2g$, reflecting the $2g$
extra $1$-cycles on $\Sigma_g$.
\end{remark}

% ----------------------------------------------------------------
\subsection{Degree $2$: the genus-$2$ KZB local system}
\label{subsec:g2-degree2}

At degree~$2$, the configuration space
$\Conf_2^{\mathrm{ord}}(\Sigma_2) = \Sigma_2^2 \setminus
\Delta$ carries the KZB connection with non-trivial
pairwise interaction. After separating the
centre-of-mass coordinate (translation on $\Sigma_2$ via
the Jacobian $\mathrm{Jac}(\Sigma_2) =
\CC^2/(\ZZ^2 + \Omega\,\ZZ^2)$, which has dimension~$2$),
the relative configuration is a punctured
genus-$2$ surface:
the KZB local system lives on $\Sigma_2 \setminus \{0\}$,
a genus-$2$ surface with one puncture.

The ordered chiral chain complex at degree~$2$ is
\begin{equation}\label{eq:g2-degree2-cx}
  \cC_2^{\mathrm{ord}}(\Sigma_2, Y_\hbar)
  = \Bigl(
  \Omega^*_{\log}\bigl(
  \overline{\FM}_2^{\mathrm{ord}}(\Sigma_2)\bigr)
  \otimes (s^{-1}\bar{Y}_\hbar)^{\otimes 2},
  \;\;
  d_{\mathrm{dR}} + d_{\mathrm{bar}}
  + \nabla_{\mathrm{KZB}}^{(g=2)}
  \Bigr).
\end{equation}
The collision residue is governed by the genus-$2$
$r$-matrix
% PE-1: r-matrix. Family: affine KM sl_2 on genus-2 curve.
%   r(z,w) = hbar * S_2(z,w) * Omega (KZ convention).
%   Level parameter: hbar = 1/(k+2).
%   Convention: KZ normalization (consistent with genus-1 section).
%   k=0 check: hbar = 1/2, r nonzero for non-abelian g (KZ: correct).
%   k=-2 check: hbar diverges (critical level: correct).
%   Source: CEE09, Bernard88.
%   verdict: ACCEPT
\begin{equation}\label{eq:r-g2}
  r^{(g=2)}(z_i, z_j)
  = \hbar\,S_2(z_i, z_j)\,\Omega
  = \frac{1}{k + h^\vee}\,
  d_{z_i}\!\log E(z_i, z_j)\,\Omega,
\end{equation}
where $\hbar = 1/(k + h^\vee) = 1/(k+2)$ and $\Omega$
is the $\mathfrak{sl}_2$ Casimir (in the KZ normalisation
consistent with equation~\eqref{eq:kzb-ell}).

\smallskip
\noindent
\textbf{Decomposition by $\mathfrak{sl}_2$-isospin.}
On $V \otimes V = \operatorname{Sym}^2 V \oplus
\bigwedge\nolimits^2 V$ with $V = \CC^2$:
\[
  \Omega\big|_{\operatorname{Sym}^2} = +\tfrac{1}{2},
  \qquad
  \Omega\big|_{\bigwedge^2} = -\tfrac{3}{2},
\]
exactly as at genus~$0$ and genus~$1$. The KZB system
decomposes into scalar equations on each isospin channel,
now with genus-$2$ coefficients.

\smallskip
\noindent
\textbf{Monodromy.}
The fundamental group
$\pi_1(\Sigma_2 \setminus \{0\})$ has presentation
\begin{equation}\label{eq:fund-gp-g2}
  \pi_1(\Sigma_2 \setminus \{0\})
  = \bigl\langle
  a_1, b_1, a_2, b_2, \gamma
  \;\big|\;
  [a_1, b_1]\,[a_2, b_2]\,\gamma = 1
  \bigr\rangle
\end{equation}
with $a_\alpha, b_\alpha$ the $A$- and $B$-cycle
generators ($\alpha = 1, 2$) and $\gamma$ the puncture
loop. The KZB connection produces a monodromy
representation
\begin{equation}\label{eq:monodromy-rep-g2}
  \rho^{(2)} \colon
  \pi_1(\Sigma_2 \setminus \{0\})
  \to \operatorname{GL}(V \otimes V)
\end{equation}
with monodromy matrices:
\begin{itemize}
\item Puncture loop:
  $M_\gamma = \exp(2\pi i\,\hbar\,\Omega)$
  (identical to genus~$0$ and genus~$1$, capturing
  the local collision residue).
\item $A$-cycles:
  $M_{a_\alpha}$ ($\alpha = 1, 2$); the $A$-cycle
  monodromy is trivial at leading order since
  $S_2(z, w)$ is periodic under $A$-cycles:
  $S_2(z + A_\alpha, w) = S_2(z, w)$.
\item $B$-cycles:
  $M_{b_\alpha}$ ($\alpha = 1, 2$); the
  quasi-periodicity~\eqref{eq:szego-quasiperiod} gives
  a shift by $-2\pi i\,\omega_\alpha$, producing
  \begin{equation}\label{eq:b-monodromy-g2}
    M_{b_\alpha}
    = \exp\!\bigl(-2\pi i\,\hbar\,
    \Omega \cdot \omega_\alpha\bigr),
    \qquad \alpha = 1, 2.
  \end{equation}
  At genus~$1$, $\omega_1 = dz$ is a scalar and
  $M_B = M_\gamma^{-1}$; at genus~$2$, the two $B$-cycle
  monodromies $M_{b_1}, M_{b_2}$ are independent and
  involve the two holomorphic differentials
  $\omega_1, \omega_2$ separately.
\end{itemize}
The constraint from the fundamental group relation is
\begin{equation}\label{eq:constraint-g2}
  [M_{a_1}, M_{b_1}]\,
  [M_{a_2}, M_{b_2}]\,
  M_\gamma = I.
\end{equation}

\begin{remark}[Dynamical parameters at genus~$2$]
\label{rem:dynamical-g2}
The $B$-cycle monodromies at genus~$2$ depend on a
\emph{two-component} dynamical parameter
$\lambda = (\lambda_1, \lambda_2) \in \fh^* \oplus \fh^*$,
replacing the single dynamical parameter $\lambda \in \fh^*$
of the genus-$1$ theory
(Remark~\textup{\ref{rem:dynamical-parameter}}).
The dynamical Yang--Baxter equation is replaced by a
system of two coupled equations, one for each $B$-cycle.
This is the genus-$2$ instantiation of the Felder dynamical
$R$-matrix theory~\cite{Felder94,FelderVarchenko96}.
\end{remark}

\begin{proposition}[Degree-$2$ ordered chiral homology on
$\Sigma_2$]
\label{prop:g2-degree2}
The degree-$2$ ordered chiral homology
$H^*(\cC_2^{\mathrm{ord}}(\Sigma_2, Y_\hbar))$ is the
de~Rham cohomology of the KZB local system on
$\Conf_2^{\mathrm{ord}}(\Sigma_2)$ with coefficients in
$(s^{-1}\bar{Y}_\hbar)^{\otimes 2}$.
On evaluation modules $V_a \otimes V_b$
\textup{(}$V = \CC^2$, so the local system has
rank~$d^2 = 4$\textup{)}:
\begin{enumerate}[label=\textup{(\roman*)}]
\item The KZB local system on
  $\Sigma_2 \setminus \{0\}$ \textup{(}a genus-$2$ surface
  with one puncture\textup{)} has monodromy
  representation~\eqref{eq:monodromy-rep-g2} with generators
  $M_{a_1}, M_{b_1}, M_{a_2}, M_{b_2}, M_\gamma$.
\item For generic $\hbar$, the monodromy has no invariant
  vectors: $H^0_{\mathrm{dR}} = 0$.
\item The Euler characteristic of a rank-$r$ local system on
  a surface of genus~$g$ with $s$ punctures is
  $\chi = r(2 - 2g - s)$. For $g = 2$, $s = 1$,
  $r = d^2 = 4$:
  \begin{equation}\label{eq:g2-degree2-euler}
    \chi = 4 \cdot (2 - 4 - 1) = -12.
  \end{equation}
\item Since $H^0 = 0$ for generic $\hbar$ and
  $H^2 = 0$ \textup{(}the puncture forces vanishing\textup{)},
  we obtain:
  \begin{equation}\label{eq:g2-degree2-h1}
    \dim H^1_{\mathrm{dR}}(\Sigma_2 \setminus \{0\},\,
    \nabla_{\mathrm{KZB}}) = 12.
  \end{equation}
\end{enumerate}
\end{proposition}

\begin{proof}
The surface $\Sigma_2 \setminus \{0\}$ has genus~$g = 2$ and
$s = 1$ punctures. The Euler characteristic of a local system
of rank~$r$ is $\chi = r(2 - 2g - s)$:
$\chi = 4(2 - 4 - 1) = 4 \cdot (-3) = -12$.
For generic $\hbar$ (equivalently, generic level~$k$), the
monodromy around the puncture is
$M_\gamma = \exp(2\pi i\,\hbar\,\Omega)$ with eigenvalues
$e^{\pi i/(k+2)}$ (on $\operatorname{Sym}^2 V$, dimension~$3$)
and $e^{-3\pi i/(k+2)}$ (on $\bigwedge^2 V$, dimension~$1$);
generically neither is~$1$, so
$\ker(M_\gamma - I) = 0$ and $H^0 = 0$.
The four $B$-cycle monodromies $M_{b_1}, M_{b_2}$ impose
further non-trivial constraints, reinforcing $H^0 = 0$.
Poincar\'e--Verdier duality on $\Sigma_2 \setminus \{0\}$
(a non-compact surface) gives
$H^2 = (H^0_c)^\vee = 0$ for generic
monodromy; alternatively, the puncture forces $H^2 = 0$ by
a residue argument.
The alternating sum gives
$0 - \dim H^1 + 0 = -12$, so $\dim H^1 = 12$.
The $12$-dimensional space decomposes along the isospin
channels:
$\operatorname{Sym}^2 V$ (rank~$3$) contributes
$3 \cdot 3 = 9$, and $\bigwedge^2 V$ (rank~$1$) contributes
$1 \cdot 3 = 3$; total $9 + 3 = 12$.
\end{proof}

\begin{remark}[Comparison of $H^1$ dimensions]
\label{rem:h1-comparison}
The degree-$2$ $H^1$ dimension grows linearly with genus:
\begin{center}
\renewcommand{\arraystretch}{1.2}
\begin{tabular}{@{} c c c c @{}}
  \toprule
  $g$ & $\chi({\Sigma_g \setminus \{0\}})$ &
  $\chi = r \cdot (2 - 2g - 1)$ & $\dim H^1$ \\
  \midrule
  $0$ & $1$ & $4 \cdot 1 = 4$ & $4$ \\
  $1$ & $-1$ & $4 \cdot (-1) = -4$ & $4$ \\
  $2$ & $-3$ & $4 \cdot (-3) = -12$ & $12$ \\
  $g$ & $1 - 2g$ & $4(1 - 2g)$ & $4(2g - 1)$ \\
  \bottomrule
\end{tabular}
\end{center}
The jump from $\dim H^1 = 4$ at $g = 1$ to $\dim H^1 = 12$
at $g = 2$ reflects the richer monodromy: at genus~$1$, the
monodromy group is generated by three elements
$(M_\gamma, M_a, M_b)$ with one relation; at genus~$2$,
by five elements $(M_\gamma, M_{a_1}, M_{b_1}, M_{a_2},
M_{b_2})$ with one relation, allowing a $12$-dimensional
space of flat sections with prescribed singular behaviour at
the puncture. The genus-$0$ entry $\dim H^1 = 4$ uses
$\chi(\PP^1 \setminus \{0, \infty\}) = 0$ for the
once-punctured sphere (the relative coordinate lives on
$\CC^\times$), so
$\chi = 0$ and the $H^1$-dimension does not follow from
$\chi$ alone; it is computed directly from KZ flat sections
\textup{(}cf.\ \S\ref{subsec:ell-degree2}\textup{)}.
\end{remark}

\begin{remark}[The symmetric shadow at genus~$2$]
\label{rem:g2-shadow}
The averaging map $\av_2$ at degree~$2$ sends the
$12$-dimensional ordered de~Rham cohomology to
$\Sigma_2$-coinvariants. The scalar surviving $\av_2$ is
$\kappa = 3(k+2)/4$
(equation~\eqref{comp:sl2-kappa}), the same value as at
genus~$0$ and genus~$1$: the averaging map is a local
computation (residue extraction), independent of the global
topology.
The kernel $\ker(\av_2)$ at genus~$2$ is richer than at
genus~$1$: the four $B$-cycle monodromies contribute
independent antisymmetric classes that are invisible to the
symmetric theory.
\end{remark}

% ----------------------------------------------------------------
\subsection{Degree $n$: genus-$2$ surface braid group
monodromy}
\label{subsec:g2-degree-n}

At degree~$n$, the fundamental group of
$\Conf_n^{\mathrm{ord}}(\Sigma_2)$ is the genus-$2$
\emph{surface braid group} $\mathrm{SB}_n(\Sigma_2)$, which
fits into the extension
\begin{equation}\label{eq:surface-braid-g2}
  1 \to PB_n \to \mathrm{SB}_n(\Sigma_2) \to
  \pi_1(\Sigma_2)^n \to 1,
\end{equation}
generalising the elliptic braid group
extension~\eqref{eq:ell-braid} (at genus~$1$,
$\pi_1(E_\tau) \cong \ZZ^2$ contributes two generators
per point; at genus~$2$,
$\pi_1(\Sigma_2) \cong \langle a_1, b_1, a_2, b_2 \mid
[a_1, b_1][a_2, b_2] = 1\rangle$ contributes four
generators per point).

The ordered chiral homology at degree~$n$ is
\begin{equation}\label{eq:g2-degree-n}
  \int_{\Sigma_2, n}^{\mathrm{ord}}
  Y_\hbar(\mathfrak{sl}_2)
  = H^*_{\mathrm{dR}}\bigl(
  \overline{\FM}_n^{\mathrm{ord}}(\Sigma_2),\;
  \nabla_{\mathrm{KZB}}^{(n)}\bigr)
  \otimes (s^{-1}\bar{Y}_\hbar)^{\otimes n}.
\end{equation}

\begin{proposition}[Euler characteristic at low degrees,
genus~$2$]
\label{prop:g2-euler-n}
The twisted Euler characteristics of the KZB local system on
$\Conf_n^{\mathrm{ord}}(\Sigma_2)$, with fibre $V^{\otimes n}$
\textup{(}$V = \CC^2$\textup{)} for generic $\hbar$, are:
\begin{align}
  n = 0&: \quad \chi = \infty
  \;\;\text{\textup{(}$\CC[\operatorname{qdet}]$ is
  infinite-dimensional\textup{)}},
  \label{eq:g2-euler-n0} \\
  n = 1&: \quad \chi = -4
  \;\;\text{\textup{(}trivial local system of rank $2$
  on $\Sigma_2$,
  $\chi(\Sigma_2) = -2$\textup{)}},
  \label{eq:g2-euler-n1} \\
  n = 2&: \quad \chi = -12
  \;\;\text{\textup{(}rank $4$ on $\Sigma_2 \setminus
  \{0\}$\textup{)}}.
  \label{eq:g2-euler-n2}
\end{align}
For $n \geq 3$, the Euler characteristic of
$\Conf_n^{\mathrm{ord}}(\Sigma_2)$ with coefficients in a
rank-$d^n$ local system is:
\begin{equation}\label{eq:g2-euler-general}
  \chi\bigl(\Conf_n^{\mathrm{ord}}(\Sigma_2),\,
  \cL_{\mathrm{KZB}}\bigr)
  = d^n \cdot \prod_{j=0}^{n-1}(\chi(\Sigma_2) - j)
  = 2^n \cdot \prod_{j=0}^{n-1}(-2 - j).
\end{equation}
In particular: $n = 3$ gives
$8 \cdot (-2)(-3)(-4) = -192$;
$n = 4$ gives $16 \cdot (-2)(-3)(-4)(-5) = -1920$.
\end{proposition}

\begin{proof}
At $n = 1$: $\Conf_1(\Sigma_2) = \Sigma_2$ and the KZB
connection is trivial (no pairwise terms).
The local system has rank~$d = 2$ and the Euler
characteristic is
$\chi = d \cdot \chi(\Sigma_2) = 2 \cdot (-2) = -4$.
(In contrast to genus~$1$ where $\chi(E_\tau) = 0$
gives $\chi = 0$, the negative Euler characteristic of
$\Sigma_2$ makes the twisted Euler characteristic non-zero
even at degree~$1$.)

At $n = 2$: Proposition~\ref{prop:g2-degree2}.

At $n \geq 3$: the product formula for the Euler
characteristic of ordered configuration spaces gives
$\chi(\Conf_n^{\mathrm{ord}}(\Sigma_2))
= \prod_{j=0}^{n-1}(\chi(\Sigma_2) - j)
= \prod_{j=0}^{n-1}(-2 - j)$, and
the multiplicativity of the Euler characteristic with
respect to local systems
gives~\eqref{eq:g2-euler-general}.
(Unlike genus~$1$, where $\chi(E_\tau) = 0$ kills the
product at the first factor and gives
$\chi(\Conf_n^{\mathrm{ord}}(E_\tau)) = 0$ for all
$n \geq 1$, here $\chi(\Sigma_2) = -2 \neq 0$ and the
product is non-vanishing for all~$n$.)
\end{proof}

% ----------------------------------------------------------------
\subsection{Summary: genus $0$ vs.\ genus $1$ vs.\
genus $2$}
\label{subsec:genus-comparison-table}

\begin{center}
\renewcommand{\arraystretch}{1.3}
\small
\begin{tabular}{@{} l l l l @{}}
\toprule
& Genus $0$ ($\PP^1$ or $D^\times$)
& Genus $1$ ($E_\tau$)
& Genus $2$ ($\Sigma_2$)
\\
\midrule
Base curve
  & $\PP^1$
  & $E_\tau = \CC/(\ZZ + \tau\ZZ)$
  & $\Sigma_2$
\\
Period matrix
  & ---
  & $\tau \in \HHH$, scalar
  & $\Omega \in \HHH_2$, $2 \times 2$
\\
Bar propagator
  & $d\log(z_1 - z_2)$
  & $d\log E(z,w) = \wp_1\,dz$
  & $d\log E(z,w)$ (prime form)
\\
\quad weight
  & $1$
  & $1$
  & $1$
\\
Connection
  & KZ
  & KZB
  & KZB (genus $2$)
\\
$\pi_1$ per point
  & $0$
  & $\ZZ^2$ ($A, B$-cycle)
  & $\ZZ^4$ ($A_1, B_1, A_2, B_2$)
\\
$B$-cycle monodromy
  & ---
  & $M_B = e^{-2\pi i\,\hbar\,\Omega}$
  & $M_{b_\alpha} = e^{-2\pi i\,\hbar\,
  \Omega \cdot \omega_\alpha}$
\\
Dynamical parameter
  & ---
  & $\lambda \in \fh^*$, one component
  & $(\lambda_1, \lambda_2) \in (\fh^*)^2$
\\
\midrule
Degree $0$
  & $\CC[\operatorname{qdet}]$
  & $\CC[\operatorname{qdet}]$
  & $\CC[\operatorname{qdet}]$
\\
Degree $1$ dim
  & $3$ (on $D^\times$)
  & $12 = 4 \times 3$
  & $18 = 6 \times 3$
\\
Degree $2$: $\dim H^1$
  & $4$
  & $4$
  & $12$
\\
Degree $2$: $\chi$
  & $4$ ($g=0,\,s=2$)
  & $-4$
  & $-12$
\\
$\chi(\Conf_n)$, $n \geq 1$
  & varies
  & $0$ (all $n \geq 1$)
  & $\prod_{j=0}^{n-1}(-2-j) \neq 0$
\\
Curvature $d_{\mathrm{fib}}^2$
  & $0$
  & $\kappa(\cA) \cdot \omega_1$
  & $\kappa(\cA) \cdot \omega_2$
\\
\quad lives on
  & ---
  & $\overline{\cM}_{1}$
  & $\overline{\cM}_{2}$
\\
Degeneration
  & ---
  & $\tau \to i\infty$: $E_\tau \to \CC^\times$
  & $\Omega_{12} \to 0$, $\Omega_{22} \to
  i\infty$: $\Sigma_2 \to E_\tau$
\\
\bottomrule
\end{tabular}
\end{center}

\begin{remark}[Non-vanishing Euler characteristics at
genus~$2$]
\label{rem:g2-nonvanishing-chi}
The qualitative difference between genus~$1$ and genus~$2$
at degree $n \geq 3$ is striking. At genus~$1$, the
vanishing $\chi(E_\tau) = 0$ forces
$\chi(\Conf_n^{\mathrm{ord}}(E_\tau)) = 0$ for all
$n \geq 1$, so the Euler characteristic provides no
information about the individual cohomology dimensions.
At genus~$2$, $\chi(\Sigma_2) = -2 \neq 0$ and the product
formula gives non-zero (and rapidly growing) Euler
characteristics: $-4, -12, -192, -1920, \ldots$ at degrees
$1, 2, 3, 4, \ldots$. These provide genuine lower bounds on
the de~Rham cohomology dimensions of the KZB local system,
making the genus-$2$ theory computationally richer: the
alternating sum is a useful invariant at every degree.
\end{remark}

\begin{remark}[Degeneration $\Sigma_2 \to E_\tau$]
\label{rem:g2-degeneration}
In the separating degeneration
$\Sigma_2 \to E_\tau \cup E_{\tau'}$ (a genus-$2$ curve
pinches along a separating cycle to two elliptic curves
joined at a node), the genus-$2$ KZB connection splits into
two independent genus-$1$ KZB connections.
The ordered chiral homology satisfies
a factorisation constraint at this boundary stratum of
$\overline{\cM}_2$:
\[
  \int_{\Sigma_2}^{\mathrm{ord}} Y_\hbar
  \;\xrightarrow{\;\mathrm{res}\;}\;
  \int_{E_\tau}^{\mathrm{ord}} Y_\hbar
  \;\otimes\;
  \int_{E_{\tau'}}^{\mathrm{ord}} Y_\hbar.
\]
In the non-separating degeneration $\Sigma_2 \to E_\tau$
(a single cycle pinches, reducing the genus by one),
the period matrix
$\Omega = \bigl(\begin{smallmatrix}
\tau & \epsilon \\ \epsilon & i/\epsilon
\end{smallmatrix}\bigr)$
degenerates as $\epsilon \to 0$, and the genus-$2$ theory
reduces to the genus-$1$ theory of
\S\ref{sec:elliptic-ordered} with additional data from
the pinched cycle.
This non-separating degeneration carries genuinely
genus-$2$ information (the off-diagonal entry
$\Omega_{12}$); the separating degeneration does not.
\end{remark}


% ================================================================
% MOVED FROM between ``Examples'' and ``Chiral quantum groups''; now after elliptic curves
\section{Comparison with related work}
\label{sec:comparison}
